\subsection{2026-09-22 --- Commitments are blocks: hour-level temporal structure}
\label{sec:temporal}

\subsubsection*{What we did}

Section~\ref{sec:abcd} concluded that the current information set was saturated and that further
architecture work was not worth doing. Looking instead for information we had never given the
model, we checked what the pointwise model can actually see at the hour level. It receives 15
engineered features, and every one of them is price-versus-water-value economics
(\texttt{f05}, \texttt{f06}, \texttt{f07}, \texttt{f10}, \texttt{f14}, \texttt{f15}). There is no
hour-of-day, no lag, no persistence, no window aggregate: each of the 168 hours is scored as an
independent draw.

\subsubsection*{Finding --- the independence assumption is wrong by an order of magnitude}

\begin{finding}
Commitments are contiguous blocks, not scattered hours.
\[
  P(\text{on at } t \mid \text{on at } t{-}1) = 0.995 \qquad
  P(\text{on at } t \mid \text{off at } t{-}1) = 0.007
\]
Units switch state \textbf{0.5--3.6 times per 168-hour week}, against 9.3--27.2 if the same hours
were scattered at random --- roughly \textbf{9.5x more contiguous than chance}. Each unit-week
contains 1.2--2.7 contiguous ON blocks.
\end{finding}

The start costs in the topology are the obvious cause: the optimiser will not pay to cycle a unit
hourly. The effective dimensionality of a unit-week is therefore not 168 bits but roughly a count
plus two or three changepoints --- and we had been modelling it as 168 independent decisions.

\subsubsection*{The obvious fix is wrong}

If commitments are blocks, allocate each unit's $c$ hours as one contiguous window. Tested
against the current rule (the $c$ highest-spread hours), both given ORACLE counts:

\begin{center}\small
\begin{tabular}{lr}
\toprule
allocation rule (oracle counts, 2021) & micro-AUC \\
\midrule
top-$c$ spread hours (current) & \textbf{0.9485} \\
best contiguous $c$-window     & 0.9057 \\
\bottomrule
\end{tabular}
\end{center}

\begin{finding}
Imposing a single contiguous window is \textbf{worse} than ranking hours by spread, by $-0.0428$.
1.2--2.7 blocks is not 1 block, and a one-window shape is too rigid. Persistence is real but does
not license a hard allocation rule.
\end{finding}

(This oracle is built differently from the 0.9765 of Section~\ref{sec:oracle} --- hard 0/1 on
2021 using system spread for every unit --- so the two figures are not comparable. Only the two
rows above are comparable to each other.)

\subsubsection*{What did work: the same structure as features}

\texttt{src/hourly\_structure.py} gives the model, per hour and always within a case: centred
window means at $w \in \{3,6,12,24,48,96\}$ with their ranks and best-window membership; run
structure (position within, and length of, the current above-median stretch); distance to the
week's peak and trough; cumulative share of the week's positive spread; lags and leads at 1, 3
and 24 hours as differences; and hour-of-day and day-of-week cyclically encoded. All are
scale-free or differenced against the case's own level, so none reintroduces the drifting price
level of Section~\ref{sec:diagnosis}.

\begin{center}\small
\begin{tabular}{lrrr}
\toprule
arm & feats & micro-AUC & Lio \\
\midrule
base (15 pruned $+$ 84 count) & 99  & 0.9417 & 0.8157 \\
$+$ hourly structure          & 135 & \textbf{0.9442} & 0.8190 \\
$+$ count-coupled windows     & 127 & 0.9418 & 0.8180 \\
$+$ both                      & 163 & 0.9430 & 0.8181 \\
\bottomrule
\end{tabular}
\end{center}

\begin{finding}
Hour-level temporal structure is worth \textbf{$+0.0025$} (0.9417 $\to$ 0.9442), the first
confirmed improvement since Section~\ref{sec:compact} and larger than A, B, C and D combined.
The gain is in the generic window and run features; \emph{coupling} the window length to the
predicted count adds nothing ($+0.0001$), so the value is in knowing that hours come in stretches,
not in knowing how long this unit's stretch will be.
\end{finding}

Per-unit: Vinje\_G3 $+0.014$, Haukeli $+0.005$, Hogga $+0.004$, Lio $+0.003$, Songa $-0.009$.

\subsubsection*{Screening harness, re-instantiated}

\texttt{src/screening.py} restores the convention of Section~\ref{sec:cv-harness}: every third
training case and 150 trees, used to rank candidates, never to score them. Strided by CASE rather
than by row, because dropping rows within a case would leave that case's near-duplicate hours in
the training set --- the leak of Section~\ref{sec:state}.

On this experiment it predicted $+0.0034$ where the full fit gave $+0.0025$: right sign, right
magnitude, 18s against 108s. Screening scores are biased low (base 0.9365 screened vs 0.9417
full) and are labelled as such wherever they appear.

\subsubsection*{Negative results from the same session}

Two families built from data identified as unused in Section~\ref{sec:diagnosis}, both rejected:

\begin{itemize}
  \item \textbf{67-year inflow climatology} (\texttt{exogenous.inflow\_climatology}, 171
        features). Day-of-year reference distributions from the full 1958--2025 record instead of
        the $\sim$8 years of case dates. Micro 0.9415 vs 0.9417 base; Lio $+0.005$.
  \item \textbf{Minimum-flow obligation} (\texttt{exogenous.minflow\_obligation}, 99 features).
        Resolved to the plant that must satisfy each constraint, converted to a forced capacity
        factor. Micro 0.9392. \textbf{Killed by arithmetic before the model saw it}: Lio at
        minimum load passes 4.82 m$^3$/s against a Byrtevatn obligation of 0.07 (max 0.20), so
        \texttt{mfo\_forced\_cf\_Lio} never exceeds 0.042 and the constraint never binds. A
        physically sound hypothesis at the wrong order of magnitude.
\end{itemize}

Also rejected: \textbf{per-unit probability calibration}. Lio is the largest single drag on the
pooled score (removing it lifts micro by $+0.0077$; an oracle Lio is worth $+0.0146$) and is
over-predicted on average (mean 0.612 against a base rate of 0.545), which suggests a logit
offset. It fails, and instructively:

\begin{center}\small
\begin{tabular}{lr}
\toprule
& micro-AUC \\
\midrule
uncalibrated & \textbf{0.9417} \\
oracle per-unit calibration (true 2021 base rates) & 0.9346 \\
legal, training base rates & 0.9363 \\
legal, Lio only & 0.9288 \\
\bottomrule
\end{tabular}
\end{center}

\begin{finding}
Even \emph{oracle} calibration hurts. A per-unit logit shift is uniform across cases, so it
cannot fix within-case ordering, and it disturbs cross-unit interleaving that was already roughly
right. A monotone per-unit transform also leaves own-AUC exactly unchanged (Lio $0.8157 \to
0.8157$). \textbf{Lio cannot be fixed by post-processing; only by better discrimination.} Same
family of error as the rank-transform failure of Section~\ref{sec:state}.
\end{finding}

\subsubsection*{What this implies}

The count stage contributes $+0.106$ (pointwise alone 0.8346 $\to$ 0.9402 with counts); everything
else combined contributes an order of magnitude less. The count model has been doing nearly all
the work, and the pointwise model has been little more than an allocator with 15 features.
Temporal structure is the first evidence that the allocation half has real headroom.

Next is the long-format restructure (\texttt{src/longform.py}): one row per (case, hour, unit),
which removes shared splits, collapses the 84-column count block to 6 \texttt{own\_} columns, and
makes per-unit weighting possible for the first time. Its main risk is losing the cross-unit
interaction a wide row carries for free, which \texttt{longform.cross\_unit\_context} attempts to
restore as relative quantities.
